% GENERATED FILE. Edit canonical lessons, then run npm run generate:books.
% canonical-source-hash: fnv1a64:f087794a83de9dd1
% canonical-lessons: PA-C76-samundar, PA-C76-agg, PA-C76-baraf, PA-C76-mimh, PA-C76-hava

\chapter{Sea, Fire, Snow, Rain, Wind}
\label{ch:pa-sea-fire-snow-rain-wind}
\hlchaptermodality{76}

\begin{chapteropening}
\textbf{By the end of this chapter:} \emph{I can say the sea, fire, snow, rain and wind, air.}

Complete the last lesson of chapter 76: I can say the sea, fire, snow, rain and wind, air.
\end{chapteropening}

\section[\texorpdfstring{\pa{ਸਮੁੰਦਰ} (samundar)}{ਸਮੁੰਦਰ (samundar)}]{\pa{ਸਮੁੰਦਰ} (samundar) — the sea}

\label{lesson:PA-C76-samundar}



\begin{quote}
Before the new one: say the Punjabi for the sun, then the Punjabi for the moon.
\end{quote}

\subsection*{You'll want to know: \pa{ਸਮੁੰਦਰ}}
\textbf{\pa{ਸਮੁੰਦਰ}} — \emph{samundar} — \textquotedblleft{}the sea\textquotedblright{}.

\begin{cousinweb}
\textbf{\pa{ਸਮੁੰਦਰ}} (\emph{samundar}) is a semi-tatsama: Sanskrit \emph{samudra}, \textquotedblleft{}sea\textquotedblright{}, borrowed after the old sound changes had run their course and then lightly reshaped (CDIAL 13236). Punjabi adds an \emph{n} and breaks up the \emph{dr}; CDIAL lists no Punjabi form.
\end{cousinweb}

\begin{grammarlens}[title={What you've built}]
One more word for this chapter.
\end{grammarlens}

\subsection*{Your turn}
\begin{itemize}
\raggedright
  \item \emph{Say it:} \emph{samundar}
  \item \emph{Say it:} \emph{samundar}, once more
  \item \emph{Say it:} say \emph{samundar}
  \item \emph{Recall:} say the Punjabi for the sun, then the Punjabi for the moon, then say \emph{samundar} again
\end{itemize}

\begin{tcolorbox}[breakable,colback=teal!4,colframe=teal!35!black,title={Before you move on}]
What does \pa{ਸਮੁੰਦਰ} mean? (\textquotedblleft{}The sea\textquotedblright{}.) Say it once more.
\end{tcolorbox}
\section[\texorpdfstring{\pa{ਅੱਗ} (agg)}{ਅੱਗ (agg)}]{\pa{ਅੱਗ} (agg) — fire}

\label{lesson:PA-C76-agg}



\begin{quote}
Before the new one: say the Punjabi for the moon, then the Punjabi for the sea.
\end{quote}

\subsection*{You'll want to know: \pa{ਅੱਗ}}
\textbf{\pa{ਅੱਗ}} — \emph{agg} — \textquotedblleft{}fire\textquotedblright{}.

\begin{cousinweb}
\textbf{\pa{ਅੱਗ}} (\emph{agg}) is a tadbhava, a word inherited from Sanskrit and reshaped by sound change. It comes from \emph{agni}, \textquotedblleft{}fire\textquotedblright{} (CDIAL 55).
\end{cousinweb}

\begin{grammarlens}[title={What you've built}]
One more word for this chapter.
\end{grammarlens}

\subsection*{Your turn}
\begin{itemize}
\raggedright
  \item \emph{Say it:} \emph{agg}
  \item \emph{Say it:} \emph{agg}, once more
  \item \emph{Say it:} say \emph{agg}
  \item \emph{Recall:} say the Punjabi for the moon, then the Punjabi for the sea, then say \emph{agg} again
\end{itemize}

\begin{tcolorbox}[breakable,colback=teal!4,colframe=teal!35!black,title={Before you move on}]
What does \pa{ਅੱਗ} mean? (\textquotedblleft{}Fire\textquotedblright{}.) Say it once more.
\end{tcolorbox}
\section[\texorpdfstring{\pa{ਬਰਫ਼} (baraf)}{ਬਰਫ਼ (baraf)}]{\pa{ਬਰਫ਼} (baraf) — snow; ice}

\label{lesson:PA-C76-baraf}



\begin{quote}
Before the new one: say the Punjabi for the sea, then the Punjabi for fire.
\end{quote}

\subsection*{You'll want to know: \pa{ਬਰਫ਼}}
\textbf{\pa{ਬਰਫ਼}} — \emph{baraf} — \textquotedblleft{}snow; ice\textquotedblright{}.

\begin{cousinweb}
\textbf{\pa{ਬਰਫ਼}} (\emph{baraf}) is a loan, usually traced to Persian \emph{barf}, \textquotedblleft{}snow, ice\textquotedblright{}.
\end{cousinweb}

\begin{grammarlens}[title={What you've built}]
One more word for this chapter.
\end{grammarlens}

\subsection*{Your turn}
\begin{itemize}
\raggedright
  \item \emph{Say it:} \emph{baraf}
  \item \emph{Say it:} \emph{baraf}, once more
  \item \emph{Say it:} say \emph{baraf}
  \item \emph{Recall:} say the Punjabi for the sea, then the Punjabi for fire, then say \emph{baraf} again
\end{itemize}

\begin{tcolorbox}[breakable,colback=teal!4,colframe=teal!35!black,title={Before you move on}]
What does \pa{ਬਰਫ਼} mean? (\textquotedblleft{}Snow; ice\textquotedblright{}.) Say it once more.
\end{tcolorbox}
\section[\texorpdfstring{\pa{ਮੀਂਹ} (mīṁh)}{ਮੀਂਹ (mīṁh)}]{\pa{ਮੀਂਹ} (mīṁh) — rain}

\label{lesson:PA-C76-mimh}



\begin{quote}
Before the new one: say the Punjabi for fire, then the Punjabi for snow; ice.
\end{quote}

\subsection*{You'll want to know: \pa{ਮੀਂਹ}}
\textbf{\pa{ਮੀਂਹ}} — \emph{mīṁh} — \textquotedblleft{}rain\textquotedblright{}.

\begin{cousinweb}
\textbf{\pa{ਮੀਂਹ}} (\emph{mīṁh}) is a tadbhava, a word inherited from Sanskrit and reshaped by sound change. It comes from \emph{megha}, \textquotedblleft{}cloud, rain\textquotedblright{} (CDIAL 10302). Punjabi kept the sense \textquotedblleft{}rain\textquotedblright{}.
\end{cousinweb}

\begin{grammarlens}[title={What you've built}]
One more word for this chapter.
\end{grammarlens}

\subsection*{Your turn}
\begin{itemize}
\raggedright
  \item \emph{Say it:} \emph{mīṁh}
  \item \emph{Say it:} \emph{mīṁh}, once more
  \item \emph{Say it:} say \emph{mīṁh}
  \item \emph{Recall:} say the Punjabi for fire, then the Punjabi for snow; ice, then say \emph{mīṁh} again
\end{itemize}

\begin{tcolorbox}[breakable,colback=teal!4,colframe=teal!35!black,title={Before you move on}]
What does \pa{ਮੀਂਹ} mean? (\textquotedblleft{}Rain\textquotedblright{}.) Say it once more.
\end{tcolorbox}
\section[\texorpdfstring{\pa{ਹਵਾ} (havā)}{ਹਵਾ (havā)}]{\pa{ਹਵਾ} (havā) — wind, air}

\label{lesson:PA-C76-hava}



\begin{quote}
Before the new one: say the Punjabi for snow; ice, then the Punjabi for rain.
\end{quote}

\subsection*{You'll want to know: \pa{ਹਵਾ}}
\textbf{\pa{ਹਵਾ}} — \emph{havā} — \textquotedblleft{}wind, air\textquotedblright{}.

\begin{cousinweb}
\textbf{\pa{ਹਵਾ}} (\emph{havā}) is a loan from Arabic \emph{hawāʾ}, \textquotedblleft{}air\textquotedblright{}, which reached Punjabi through Persian. Its Arabic root is \emph{h-w-y} (Lane).
\end{cousinweb}

\begin{grammarlens}[title={What you've built}]
One more word for this chapter.
\end{grammarlens}

\subsection*{Your turn}
\begin{itemize}
\raggedright
  \item \emph{Say it:} \emph{havā}
  \item \emph{Say it:} \emph{havā}, once more
  \item \emph{Say it:} say \emph{havā}
  \item \emph{Recall:} say the Punjabi for snow; ice, then the Punjabi for rain, then say \emph{havā} again
\end{itemize}

\begin{tcolorbox}[breakable,colback=teal!4,colframe=teal!35!black,title={Before you move on}]
What does \pa{ਹਵਾ} mean? (\textquotedblleft{}Wind, air\textquotedblright{}.) Say it once more.
\end{tcolorbox}
